% Generated draft: requires mathematical and source review.
\documentclass[11pt]{article}
\usepackage[margin=25mm]{geometry}
\usepackage{fontspec}
\setmainfont{Latin Modern Roman}
\usepackage{amsmath,amssymb,mathtools,unicode-math}
\setmathfont{Latin Modern Math}
\usepackage{enumitem,needspace,xurl,hyperref}
\setlength{\parindent}{0pt}
\setlength{\parskip}{5pt}
\setlength{\emergencystretch}{3em}
\newcommand{\N}{\mathbb{N}}
\newcommand{\Z}{\mathbb{Z}}
\newcommand{\Q}{\mathbb{Q}}
\newcommand{\R}{\mathbb{R}}
\newcommand{\C}{\mathbb{C}}
\newcommand{\F}{\mathbb{F}}
\newcommand{\A}{\mathbb{A}}
\newcommand{\PP}{\mathbb{P}}
\newcommand{\E}{\mathbb{E}}
\newcommand{\eps}{\varepsilon}
\newcommand{\dd}{\,\mathrm{d}}
\begin{document}
\section*{UnsolvedMath ID: 30003013}
\textbf{Source problem code:} catalogue entry
\subsection{Definitions and mathematical statement}
Let
\[
\mathfrak H=\{\tau\in\C:\operatorname{Im}\tau>0\}
\]
be the complex upper half-plane. For an algebraic modular curve \(X\) and a line bundle \(\omega\) on \(X\), the notation \(H^0(X,\omega^{\otimes k})\) denotes the space of global sections of the indicated tensor power. In the usual modular-curve interpretation, \(\omega\) is the Hodge bundle, whose fiber at a point representing an elliptic curve \(E\) is the space of invariant differentials on \(E\). This interpretation requires a suitable moduli space or stack; the supplied record does not specify one. The symbol \(k\) denotes the weight in the source, whose allowed values are not stated.

\textbf{Source question.} For a holomorphic modular form, explain how its description as an algebraic global section
\[
s\in H^0(X,\omega^{\otimes k})
\]
and its description as a holomorphic modular function \(f:\mathfrak H\to\C\) of weight \(k\) can be identified or related. Here ``modular function'' means the function representing the modular form, and does not impose weight zero.

\textbf{Unresolved qualifications.} The source specifies neither the level or modular group, the algebraic base field and complex realization, the moduli interpretation of \(X\), nor whether \(X\) is compactified. It also omits the transformation law, any character or multiplier, the conditions at cusps, the extension of \(\omega\) over a compactification, and the normalization used to express a section as a function. Consequently, the displayed descriptions identify the objects to be compared but do not determine a precise comparison map or a proposition asserting that such a map is an isomorphism. The supplied question is therefore retained as a request for a comparison, with its missing hypotheses unresolved.
\subsection{Short English statement}
The question asks how the algebraic-section and upper-half-plane descriptions of a holomorphic modular form correspond. The supplied record leaves the level, weight conventions, geometric model, and cusp conditions unspecified, so it does not determine a single precise comparison assertion.
\subsection{Sources}
\begin{itemize}
\item Supplied local UnsolvedMath registry record, permanent ID \(30003013\), title \emph{Comparison of Algebraic and Analytic Modular Forms}, source problem code \texttt{catalogue entry}, in the Oberwolfach Reports Open Problems collection. No individual proposer is supplied; the proposed year is recorded as 2015.
\item Supplied report citation: \emph{Non-Archimedean Geometry and Applications} (2016), Oberwolfach Reports, DOI 10.4171/owr/2015/57, pp.\ 3271--3331. Supplied source URL and literature-source link: \url{https://doi.org/10.4171/owr/2015/57}. The background describes the report as documenting mathematical workshops and their open-problem discussions; no passage or problem-specific page is supplied.
\item Literature assessment dated exactly \texttt{2026-08-21}: the supplied curation labels the item open because the extracted text is a broad transition question between viewpoints and does not specify an assertion whose truth can be checked. This is supplied dated curation, not an independently verified literature assessment. The supplied background requests independent verification of time-sensitive status and cited claims.
\item The registry labels the statement \texttt{corrected\_verified} and says that mathematical content was preserved while notation, typography, or omitted definitions were repaired. These are supplied editorial labels and claims, not mathematical verification of this draft or its intended comparison.
\end{itemize}
\end{document}
